\chapter{Latex使用}

\section{原样输出}
\begin{minted}{tex}
   行内原样输出:\par
    \verb|\verb||代码或文字|

    多行原样输出：\par
    \begin{lstlisting}
    \begin{verbatim}
    多行代码
    \end{verbatim}
    \end{lstlisting}

 
\end{minted}

\section{代码显示}
可以使用宏包\verb|\usepackage{listing}|或者\verb|\usepackage{minted}|\par
但是添加宏包后，需要对宏包进行参数设置，添加配置，本文档使用了minted宏包并在note.sty中进行了配置
\section{字体颜色}

\begin{verbatim}
\textcolor{red}{红色文字}
\end{verbatim}


\section{注释}
\verb|\footnote{正文部分}| 

\section{网址链接}
\verb|\url{https://www.example.com}|\par
\textcolor{red}{对于较长的网址，要保证正文对其，在配置页面添加新的package}，\verb|\usepackage{xurl}|

\section{超链接}
\verb|\href{https://www.example.com}{点击这里访问网站}|\par
\textcolor{blue}{超链接，不显示网址，直接点击文字跳转}

\section{自定义颜色}

\section{图片导入}
\begin{minted}{tex}
    \begin{figure}[htbp]
    \centering
    \includegraphics[width=0.9\linewidth]{figure/sw_load.png}  //[]中的是反斜杠
    \caption{设置雷达芯片型号}
    \label{fig:设置雷达芯片型号}
    \end{figure}
\end{minted}


\section{表格}
\begin{minted}{tex}
  \begin{table}[htbp]
    \centering
    begin{tabular}{|c|c|c|}
    \hline
    模式&  描述&开关设置\\
    \hline
    1.刷写模式& 雷达芯片初始化-USB &S1: 0100  S4: 0000\\
    \hline
    \end{tabular}
    \caption{刷写模式开关设置}
    \label{tab:刷写模式开关设置}
\end{table}  
\end{minted}

